\subsection{Golden Section Search}

% ---------------------------------------------------------------
% 1. Explicación gráfica
% ---------------------------------------------------------------
\begin{frame}{Golden Section Search: Explicación Gráfica}
    \begin{block}{Concepto principal}
        \small
        Como $F_n/F_{n-1}\to\varphi$, se reemplaza la razón variable de Fibonacci por la
        \textbf{razón áurea fija} $\rho=\varphi-1\approx 0.618$. Al ser la misma en toda
        iteración, la reducción es autosemejante: el punto interior que sobrevive queda
        ya colocado en la posición áurea del nuevo intervalo y \textbf{se recicla}.
    \end{block}

    \begin{center}
    \begin{tikzpicture}[x=0.82cm, y=1cm, font=\scriptsize]
        % Intervalo I_k
        \draw[CelesteBrillante, thick] (0,0) -- (8,0);
        \foreach \x/\lb in {0/a, 3.056/c, 4.944/d, 8/b}{
            \draw[CelesteBrillante, thick] (\x,-0.13) -- (\x,0.13);
            \node[above, text=GrisTexto] at (\x,0.13) {$\lb$};
        }
        \node[right, text=GrisTexto] at (8.25,0) {$I_k$};
        % Reutilización del punto c
        \draw[NaranjaBase, dashed, thick, ->] (3.056,-0.25) -- (3.056,-1.05);
        \node[right, text=NaranjaBase] at (3.3,-0.65) {$c$ se reutiliza};
        % Nuevo intervalo tras f(c) < f(d)
        \draw[CelesteBrillante, thick] (0,-1.3) -- (4.944,-1.3);
        \foreach \x/\lb in {0/a, 1.888/c, 3.056/d, 4.944/b}{
            \draw[CelesteBrillante, thick] (\x,-1.43) -- (\x,-1.17);
            \node[below, text=GrisTexto] at (\x,-1.43) {$\lb$};
        }
        \node[right, text=GrisTexto] at (5.2,-1.3) {$I_{k+1}=I_k/\varphi$};
    \end{tikzpicture}
    \end{center}

    \vspace{0.2em}
    \footnotesize
    Solo el nuevo $c$ requiere evaluar $f$; las razones $0.382$ y $0.618$ se repiten
    idénticas en cada paso.
\end{frame}

% ---------------------------------------------------------------
% 2. Pseudocódigo
% ---------------------------------------------------------------
\begin{frame}{Golden Section Search: Pseudocódigo}
    \small
    \begin{algorithmic}[1]
        \Require $f$ unimodal, intervalo $[a,b]$, número de evaluaciones $n>1$
        \State $\rho \gets \varphi - 1 = (\sqrt{5}-1)/2 \approx 0.618$
        \State $d \gets \rho\, b + (1-\rho)\, a$; \quad $y_d \gets f(d)$
              \Comment{única evaluación previa}
        \For{$i = 1$ \textbf{to} $n-1$}
            \State $c \gets \rho\, a + (1-\rho)\, b$ \Comment{punto izquierdo}
            \State $y_c \gets f(c)$
            \If{$y_c < y_d$}
                \State $(b,\, d,\, y_d) \gets (d,\, c,\, y_c)$ \Comment{el mínimo está del lado de $c$}
            \Else
                \State $(a,\, b) \gets (b,\, c)$ \Comment{se invierte la orientación}
            \EndIf
        \EndFor
        \State \Return $[\min(a,b),\ \max(a,b)]$
    \end{algorithmic}
    \vspace{0.4em}
    {\footnotesize \textit{Nota:} $\rho$ ya no depende de $n$; el bucle puede detenerse
     en cualquier momento o cambiarse por el criterio $|b-a|<\epsilon$.}
\end{frame}

% ---------------------------------------------------------------
% 3. Ejemplo paso a paso
% ---------------------------------------------------------------
\begin{frame}{Golden Section Search: Ejemplo Paso a Paso}
    \begin{exampleblock}{Planteamiento}
        \small
        La misma función que en Fibonacci: minimizar $f(x)=e^{x-2}-x$ en $[-2,6]$
        con $n=5$ evaluaciones, para poder comparar ambos métodos.
    \end{exampleblock}

    \vspace{-0.2em}
    \begin{center}
    \scriptsize
    \begin{tabular}{c c r c l r}
        \toprule
        \# & consulta & $f$ & comparación & nuevo $[a,b]$ & long. \\
        \midrule
        0 & $d=2.944$ & $-0.373$ & ---          & $[-2,\,6]$          & $8.000$ \\
        1 & $c=1.056$ & $-0.667$ & $y_c<y_d$    & $[-2,\,2.944]$      & $4.944$ \\
        2 & $c=-0.112$ & $+0.233$ & $y_c\ge y_d$ & $[-0.112,\,2.944]$ & $3.056$ \\
        3 & $c=1.777$ & $-0.977$ & $y_c<y_d$    & $[1.056,\,2.944]$   & $1.889$ \\
        4 & $c=2.223$ & $-0.973$ & $y_c\ge y_d$ & $[1.056,\,2.223]$   & $1.167$ \\
        \bottomrule
    \end{tabular}
    \end{center}

    \footnotesize
    Cada longitud es la anterior dividida entre $\varphi$; la final es
    $8/\varphi^{4}=1.167$, algo mayor que el $1.01$ de Fibonacci con el mismo
    presupuesto, y también contiene a $x^{*}=2$.
\end{frame}

% ---------------------------------------------------------------
% 4. Código
% ---------------------------------------------------------------
\begin{frame}[fragile]{Golden Section Search: Código}
\begin{lstlisting}[language=Python]
from math import sqrt
PHI = (1 + sqrt(5)) / 2

def golden_section_search(f, a, b, n):
    rho = PHI - 1                  # 0.6180339887...
    d = rho * b + (1 - rho) * a    # punto derecho
    yd = f(d)
    for _ in range(n - 1):
        c = rho * a + (1 - rho) * b   # punto izquierdo
        yc = f(c)
        if yc < yd:
            b, d, yd = d, c, yc   # el minimo esta junto a c
        else:
            a, b = b, c           # se invierte el intervalo
    return (a, b) if a < b else (b, a)

# Por tolerancia: n = ceil(log((b - a) / eps) / log(PHI))
# golden_section_search(lambda x: exp(x-2) - x, -2, 6, 5)
#   -> (1.0557, 2.2229)
\end{lstlisting}
\end{frame}

% ---------------------------------------------------------------
% 5. Análisis de convergencia
% ---------------------------------------------------------------
\begin{frame}{Golden Section Search: Análisis de Convergencia}
    \begin{alertblock}{Tasa de convergencia}
        \small
        Tras $n$ evaluaciones el intervalo mide
        $\;|b_n-a_n| = (b-a)\,\varphi^{-(n-1)}$.
        La razón de reducción es \textbf{constante e igual a} $1/\varphi\approx 0.618$:
        convergencia \textbf{lineal} de orden $q=1$.
    \end{alertblock}

    \footnotesize
    \begin{itemize}
        \item \textbf{Iteraciones para una tolerancia $\epsilon$:}
              $n=\left\lceil \ln\!\big((b-a)/\epsilon\big)/\ln\varphi \right\rceil$;
              p.\,ej. reducir $[-2,6]$ a $10^{-4}$ requiere $n=24$.
        \item \textbf{Frente a Fibonacci:} $\varphi^{\,n-1}$ contra $F_{n+1}$
              ($6.85$ contra $8$ si $n=5$). Fibonacci es óptimo, pero su ventaja se
              estabiliza en el factor constante
              $F_{n+1}/\varphi^{\,n-1}\to\varphi^{2}/\sqrt{5}\approx 1.17$.
        \item \textbf{Ventaja práctica:} $\rho$ es fijo, así que no hace falta conocer $n$
              de antemano y se puede parar con un criterio sobre $|b-a|$.
        \item \textbf{Limitación:} igual que Fibonacci, si $f$ no es unimodal el
              intervalo final puede encerrar solo un mínimo local.
    \end{itemize}
\end{frame}
